\section{Resultados}
\label{sec:resultados}

\subsection{Realce, ruido, frecuencia y wavelets (OE2)}
\label{sec:res-oe2}

\paragraph{Estrategias de realce y filtrado.} De las \FinRealceTotal{} estrategias
evaluadas sobre la muestra controlada, \textbf{solo \FinRealceConservadas{} cumple el
criterio de la formulación}: el pasa-bajas gaussiano con corte 0,15 ciclos/píxel,
por un margen mínimo sobre el umbral de conservación de bordes. Los filtros de
suavizado eliminan casi todo el ruido estimado a costa de los bordes; el realce
(\emph{unsharp}, ecualización, CLAHE) conserva los bordes pero amplifica el ruido.
La mediana 5$\times$5 fue la que mejor equilibró ruido y bordes, y quedó justo por
debajo del umbral. El ruido de las imágenes originales es menor que un nivel de
gris en 8 bits: hay poco ruido que quitar.

\input{final_oe2_decision}

\begin{figure}[H]
\centering
\includegraphics[width=\textwidth]{bordes_comparacion.png}
\caption{Operadores de borde (Sobel, Prewitt, Roberts y Laplaciano) sobre una
imagen de cada origen.}
\end{figure}

\paragraph{Firma espectral de cada origen.} La fracción de potencia por encima de
0,25 ciclos/píxel tiene una mediana de \FinAltaMC{} en Montgomery y \FinAltaSH{} en
Shenzhen. El delta de Cliff entre orígenes es \FinDeltaEspectro{}: \textbf{separación
total}, todas las imágenes de Montgomery tienen menos alta frecuencia que todas las
de Shenzhen. No se puede separar del todo el efecto del detector del de la
reducción de resolución, porque Shenzhen no documenta el tamaño de píxel.

\begin{figure}[H]
\centering
\includegraphics[width=\textwidth]{espectros.png}
\caption{Espectro medio de cada origen sobre un recorte central dentro del campo
(izquierda y centro) y perfil radial de potencia normalizada (derecha).}
\end{figure}

\paragraph{Energía por subbanda.} La descomposición db2 de tres niveles localiza la
diferencia: las subbandas \textbf{diagonales} se separan por completo en las tres
escalas y las \textbf{verticales} de forma marcada, mientras que las
\textbf{horizontales no difieren}. En una radiografía PA, la energía horizontal
está dominada por las costillas y el diafragma, la misma anatomía en ambos
orígenes; las orientaciones vertical y diagonal reflejan sobre todo textura y
adquisición. La firma del sitio reside en la adquisición, no en la anatomía. Como
la reducción de resolución afectaría todas las orientaciones por igual, el
resultado sugiere que la diferencia no se explica solo por ella.

\input{final_oe2_wavelets}

\begin{figure}[H]
\centering
\includegraphics[width=0.8\textwidth]{energia_wavelets.png}
\caption{Energía relativa por escala y por origen (DWT db2, tres niveles).}
\end{figure}

\paragraph{Especificación de histograma.} Llevar la distribución de intensidad de
Montgomery a la de Shenzhen elimina la diferencia de intensidad (delta de
\FinIntensidadAntes{} a \FinIntensidadDespues{}), pero la alta frecuencia sigue
separada (delta \FinAltaDespues{}). Se había predicho que la textura no cambiaría,
por ser una transformación puntual; en realidad cambió en parte, porque la
transformación no es lineal y estira el contraste local en el rango del
parénquima. \textbf{El efecto de sitio no se corrige normalizando intensidades.}

\input{final_oe2_especificacion}

\begin{figure}[H]
\centering
\includegraphics[width=\textwidth]{histogram_matching.png}
\caption{Montgomery antes y después de especificar su histograma hacia Shenzhen.}
\end{figure}

\paragraph{Geometría.} En las \FinConectN{} imágenes de la muestra, el número de
componentes de la máscara coincidió con vecindad 4 y 8 en \FinConectIguales{}
casos: la regla de elegibilidad no depende de esa elección.

\subsection{Segmentación contra la referencia (OE3)}
\label{sec:res-oe3}

En entrenamiento, \texttt{base\_adap} obtuvo el mayor Dice mediano en ambos
orígenes, con valores casi iguales entre sitios. Ninguna estrategia con realce
mejoró la línea base de forma consistente en los dos orígenes.

\input{dice}

\begin{figure}[H]
\centering
\includegraphics[width=\textwidth]{dice_por_estrategia.png}
\caption{Concordancia por estrategia en entrenamiento. En verde, la línea base.}
\end{figure}

\paragraph{Criterio de éxito.} La formulación exigía al menos una combinación cuyo
error relativo mediano mejorara la línea base en tres de los cinco descriptores,
en ambos orígenes. Ninguna lo cumple; las dos que se acercan fallan por
inconsistencia entre sitios.

\begin{table}[H]
\centering
\small
\caption{Evaluación del criterio de éxito. Cada estrategia con realce se compara
con la línea base que comparte umbralización.}
\begin{tabular}{@{}llccl@{}}
\toprule
\textbf{Estrategia} & \textbf{Línea base} & \textbf{Montgomery} & \textbf{Shenzhen} & \textbf{Resultado} \\
\midrule
clahe\_adap & base\_adap & 0/5 & 2/5 & no cumple \\
clahe\_fijo & base\_fijo & 2/5 & 4/5 & no cumple \\
clahe\_otsu & base\_otsu & 1/5 & 2/5 & no cumple \\
eq\_adap & base\_adap & 3/5 & 1/5 & no cumple \\
eq\_fijo & base\_fijo & 0/5 & 2/5 & no cumple \\
eq\_otsu & base\_otsu & 2/5 & 3/5 & no cumple \\
\bottomrule
\end{tabular}
\end{table}

Es un resultado negativo y sustentado. Con un promedio agregado en lugar de exigir
consistencia en ambos orígenes, \texttt{clahe\_fijo} habría parecido una mejora.

\paragraph{Verificación en validación.} En las imágenes de validación,
\texttt{base\_adap} mantuvo el primer puesto en ambos orígenes, con Dice mediano de
\FinDiceValMC{} en Montgomery y \FinDiceValSH{} en Shenzhen (\FinDiceTrainMC{} y
\FinDiceTrainSH{} en entrenamiento). El ranking de las nueve estrategias coincidió
con el de entrenamiento salvo un intercambio, de modo que la elección no dependió
de la muestra. En las comparaciones pareadas fue significativamente superior en
\FinParesSig{} de \FinParesTotal{} casos; no difirió significativamente de
\texttt{clahe\_fijo} en ninguno de los orígenes ni de \texttt{eq\_adap} en
Montgomery. Solo \texttt{base\_adap} es consistente entre sitios: \texttt{eq\_adap}
queda segunda en Montgomery y sexta en Shenzhen.

\input{final_oe3_ranking}
\input{final_oe3_pareadas}

\begin{figure}[H]
\centering
\includegraphics[width=\textwidth]{oe3_dice_validacion.png}
\caption{Dice por estrategia en validación. En rojo, \texttt{base\_adap}.}
\end{figure}

\subsection{Desplazamiento de los descriptores (OE4)}
\label{sec:res-oe4}

El preprocesamiento desplaza los descriptores de tres formas distintas, con el
mismo comportamiento en entrenamiento, validación y prueba:

\begin{itemize}
    \item \textbf{Sesgo sistemático.} La compacidad automática es menor que la de
    referencia en prácticamente todas las imágenes, cerca de un \FinSesgoCompacidad\,\%
    menor. El umbral deja bordes dentados donde el trazado humano es suave, y un
    borde dentado tiene muchos más píxeles de perímetro. La fracción de área y la
    energía wavelet se sobreestiman en la mayoría de las imágenes.
    \item \textbf{Sesgo dependiente del sitio.} La relación de aspecto se sobreestima
    en Shenzhen y casi no en Montgomery. Es la inestabilidad más grave para la
    pregunta, que pide mediciones reproducibles entre sitios.
    \item \textbf{Error aleatorio.} La simetría no tiene dirección (cerca de la mitad
    de las imágenes sobreestimadas), pero su error es del orden de la propia medida.
    Su error relativo explota cuando la referencia es casi cero, por lo que se
    reporta la diferencia absoluta.
\end{itemize}

\input{final_oe4}

\begin{figure}[H]
\centering
\includegraphics[width=\textwidth]{error_descriptores_particiones.png}
\caption{Error relativo con signo de cada descriptor de \texttt{base\_adap} por
partición y origen. Positivo: la segmentación automática sobreestima.}
\end{figure}

Por estrategia, en entrenamiento, la relación de aspecto es el descriptor más
estable y la simetría el menos estable:

\input{errores}

\subsection{Diferencias entre grupos (OE5)}
\label{sec:res-oe5}

Se analizaron \FinPacientesOECinco{} pacientes. Con la máscara de referencia, en
Shenzhen \textbf{\FinOECincoSigSH{} descriptores difieren} entre radiografías con y
sin hallazgos, con efectos moderados: el campo pulmonar con hallazgos es más
alargado, ocupa más fracción del campo, es menos compacto y tiene menos energía de
detalle fino. En Montgomery ninguno es significativo tras la corrección. Son
asociaciones morfométricas, no criterios diagnósticos.

\input{final_oe5_referencia}

\begin{figure}[H]
\centering
\includegraphics[width=\textwidth]{oe5_delta_cliff.png}
\caption{Delta de Cliff entre grupos con intervalo bootstrap al 95\,\%, con la
máscara de referencia y con la segmentación automática.}
\end{figure}

\paragraph{Confusión.} En Shenzhen los grupos son comparables en edad y sexo; en
Montgomery el grupo con hallazgos es mayor y tiene más hombres, de modo que
cualquier diferencia allí estaría confundida, además de tener poca potencia.

\input{final_oe5_covariables}

\begin{figure}[H]
\centering
\includegraphics[width=\textwidth]{distribuciones.png}
\caption{Distribución de los descriptores de referencia por grupo y origen, un
punto por paciente.}
\end{figure}

\paragraph{Sensibilidad a la segmentación.} Con los descriptores automáticos,
\FinEstables{} de \FinEstabilidadTotal{} conclusiones se mantienen. La compacidad
ilustra el caso más claro: su sesgo es enorme, pero su diferencia entre grupos es
casi idéntica con ambas máscaras, porque \textbf{un sesgo uniforme no altera la
comparación entre grupos}. En cambio, la relación de aspecto pierde en Shenzhen una
diferencia presente en la referencia, y en Montgomery aparece una diferencia
significativa que la referencia no muestra: su error no es igual en los dos grupos.

\input{final_oe5_estabilidad}

Con las nueve estrategias, cada una cambia entre \FinDiscMin{} y \FinDiscMax{} de
las 10 conclusiones respecto a la referencia; \texttt{base\_adap} cambia
\FinDiscBaseAdap{}. \textbf{Ninguna estrategia es inocua}, y la de mejor Dice no es
la que mejor preserva las conclusiones. En Shenzhen la segmentación tiende a borrar
diferencias reales; en Montgomery, a crear diferencias que la referencia no
muestra.

\input{final_oe5_estrategias}

\begin{figure}[H]
\centering
\includegraphics[width=\textwidth]{oe5_estrategias_train.png}
\caption{Delta de Cliff por estrategia, descriptor y origen en entrenamiento.
Asterisco: significativo tras Holm. Recuadro: no coincide con la referencia.}
\end{figure}

\subsection{Modelado supervisado}
\label{sec:res-modelo}

En validación, el modelo de grado 2 obtuvo el menor RMSE (\ModRMSEValDos{}), frente
a \ModRMSEValUno{} del grado 1 y \ModRMSEValBase{} de la media de entrenamiento. Su
ventaja sobre el grado 1 es modesta. La matriz de diseño tiene rango completo
(\ModRangoDiseno{}) y número de condición \ModCondicion{}: los coeficientes son
identificables.

\input{lab005_validacion}

En la prueba, evaluada una sola vez con el modelo congelado, el RMSE fue
\FinRMSE{} (IC 95\,\% \FinRMSEInf{} a \FinRMSESup{}), frente a \FinRMSEMedia{} de la
media de entrenamiento: una reducción del \FinReduccion\,\% (\FinReduccionInf{} a
\FinReduccionSup\,\%), con $R^2 = \ModRTest{}$. El intervalo de la reducción no
incluye el cero, y el desempeño coincide con el de validación. El error en
Montgomery es más alto, pero su intervalo, amplio por tener \ModNTestMC{} pacientes,
se superpone con el de Shenzhen.

\input{final_bootstrap}

\begin{figure}[H]
\centering
\includegraphics[width=\textwidth]{prueba_prediccion_residuos.png}
\caption{Prueba: referencia frente a predicción y residuos. Las fracciones
extremas regresan a la media.}
\end{figure}

\input{lab005_coeficientes}

\paragraph{Práctica común.} Con datos sintéticos, el modelo de \DidModeloReg{}
recuperó el término cuadrático del mecanismo (RMSE de prueba \DidRMSETestReg{}
frente a \DidRMSETestBase{} de la referencia). La regresión logística alcanzó un
ROC-AUC de \DidAUCTest{} y un AP de \DidAPTest{} frente a una prevalencia de
\DidPrevTest{}, pero con umbral 0,5 su sensibilidad fue \DidSensTest{}: el umbral
es una decisión de costos que debe fijarse en validación.

\input{lab005_didactico_reg}
\input{lab005_didactico_bin}